\documentclass{article}
\usepackage{amsmath}
\usepackage{amsfonts}

\title{SDF JAVA PROJECT}
\author{H VASANT KUMAR}
\date{}

\begin{document}

\maketitle

\section*{Objective}
\begin{itemize}
    \item Implementation of an infinite/arbitrary-precision arithmetic library in Java using OOP concepts.
    \item If you are given two strings as input, you need to add support for addition, subtraction, multiplication, and division, for both integer and float data types.
    \item Ensure efficient memory management when dealing with large integers and floating-point numbers, especially when performing operations on large datasets or long-precision numbers.
    \item  Implement a mechanism to define the precision of the arithmetic operations.
    \item Safely handle very large numbers and overflow scenarios.
\end{itemize}

\section*{Library Specification}
It has a package called \texttt{arbitraryarithmetic} which consists of two classes \texttt{AInteger} and \texttt{AFloat}.

\section*{AInteger.java / Integer class}
I created a directory named \texttt{arbitraryarithmetic}. Inside this directory, I created the file \texttt{AInteger.java}. To make it a package, I included the statement \texttt{package arbitraryarithmetic;} at the top of the file.

\subsection*{Properties or members}
It has a protected member named \texttt{'num'} of type String.

\subsection*{Constructors}
\begin{itemize}
    \item Default constructor \texttt{AInteger()} that initializes the instance with value 0.
    \item Constructor \texttt{AInteger(String num)} that initializes the instance by the number whose string representation is given by 'num'.
    \item Copy constructor that creates an instance of \texttt{AInteger}.
\end{itemize}

\subsection*{Methods / Functions}
\begin{description}
    \item[\texttt{parse(String num)}] A static function that returns an instance of \texttt{AInteger} class.
    \item[\texttt{get()}] Returns the object's member which is "num".
    \item[\texttt{set(String num)}] Sets the object's member which is "num".
    \item[\texttt{remove\_leading\_zero(String num)}] Removes the extra unnecessary leading zeroes.
    \item[\texttt{add(AInteger other)}] Adds integers of arbitrary length.
    \item[\texttt{sub(AInteger other)}] Subtracts integers of arbitrary length.
    \item[\texttt{mul(AInteger other)}] Multiplies integers of arbitrary length.
    \item[\texttt{div(AInteger other)}] Divides integers of arbitrary length.
\end{description}

\subsection*{Deeper into the methods}
\subsubsection*{\texttt{add(AInteger other)}}
\begin{itemize}
    \item I first eliminated leading zeros from both operands to ensure accurate numerical representation and comparison.
    \item This function adds any two integers, whether positive or negative.
    \item I initially ensured it worked correctly for positive integers.
    \item I Implemented manual addition by iterating through digits from right to left (least to most significant), simulating arithmetic by hand.
    \item The numbers are processed in a while loop, where each digit is added starting from the unit's place.
    \item Manages and propagates carry values accurately during each step of the digit-wise addition process.
    \item Accommodates operands of different lengths by continuing the addition with remaining digits of the longer operand once the shorter is exhausted.
    \item Accumulates the result in reverse order and subsequently reverses it to obtain the final correctly ordered result.
    \item Handles any leftover carry after completing digit-wise addition by appending it to the final result string.
    \item Prepends a negative sign to the result if both original operands were negative, preserving mathematical correctness.
    \item It initially detects scenarios where the operands have opposing signs and transforms the operation into a subtraction of absolute values.
    \item Finally, I passed the result to a function that removes any leading zeros.
    \item Returns a new \texttt{AInteger} object encapsulating the final computed result, conforming to the class's design contract.
    \item Supports addition of arbitrarily large integers by performing all operations at the string level, circumventing native data type limitations.
\end{itemize}

\subsubsection*{\texttt{sub(AInteger other)}}
\begin{itemize}
    \item First, removes any leading zeroes from both numbers to make sure the values are clean and accurate.
    \item As with addition, I initially implemented subtraction for positive integers.
    \item If both numbers are exactly the same, the result is zero.
    \item I compared the inputs to determine whether the result should be positive or negative.
    \item Stores a flag to keep track of whether the result should have a minus sign or not
    \item The subtraction was performed by subtracting the smaller number from the larger one.
    \item Performs subtraction digit by digit starting from the rightmost digit (like in manual subtraction)
    \item If a digit in the first number is smaller than the digit being subtracted, it borrows from the next digit.
    \item Keeps subtracting digits until all digits in both numbers are used up.
    \item Reverses the result string since digits were added from right to left during subtraction. 
    \item Depending on the earlier comparison, I prefixed the result with a minus sign if necessary.
    \item If both inputs were negative, I simply reversed the order of the operands before performing the subtraction.
    \item If the first number is negative and the second is positive, it adds the magnitudes and puts a minus sign in front of the result.
    \item If the first number is positive and the second is negative, it simply adds the two magnitudes (because subtracting a negative is like adding).
    \item Finally, I passed the result to a function that removes any leading zeros.
    \item Creates and returns a new \texttt{AInteger} object with the final result.
\end{itemize}

\subsubsection*{\texttt{mul(AInteger other)}}
\begin{itemize}
    \item Removes any leading zeroes from both input numbers to clean them up.
    \item If either number is zero, the result is immediately zero.
    \item I initially implemented multiplication for positive integers, building upon the logic used for addition and subtraction.
    \item Initializes an empty result that will hold the final answer.
    \item Loops through each digit of the second number from left to right.
    \item For each digit
    \begin{itemize}
        \item Converts the character digit to an integer
        \item Creates a temporary result by adding the first number to itself 'digit' number of times.
        \item Appends the appropriate number of zeroes at the end to simulate multiplication by powers of 10.
    \end{itemize}
    \item Adds the result of each digit’s multiplication to the running total.
    \item If both input numbers were negative, I multiplied their absolute values.
    \item If only one of the inputs was negative, I multiplied their absolute values and prefixed the result with a minus sign.
    \item Finally, the result was passed through a function to remove any leading zeros before returning it.
    \item Returns the result wrapped inside a new \texttt{AInteger} object.
\end{itemize}

\subsubsection*{\texttt{div(AInteger other)}}
\begin{itemize}
    \item As with other arithmetic operations, I first removed any leading zeros from the inputs.
    \item Then, I removed the negative signs so it can just work with magnitudes.
    \item I checked if the divisor was zero and, if so, threw an error to handle division by zero.
    \item If the dividend was zero, I immediately returned zero as the result.
    \item If divisor is 1, the result is just the dividend.
    \item If both are equal, the result is 1.
    \item If divisor is larger than dividend, result is 0.
    \item I then implemented the long division algorithm for performing the division.
    \item I take one digit at a time from the dividend:
    \item The code builds a string \texttt{comp} to represent a part of the dividend being divided.
    \item For each digit in the dividend:
        \begin{itemize}
            \item Adds the digit to \texttt{comp}.
            \item If \texttt{comp} is still smaller than the divisor, adds "0" to the result and continues.
            \item Otherwise, finds the largest digit (0 to 9) such that $ \text{divisor} \times \text{digit} \leq \texttt{comp} $
            \item Adds that digit to the final answer string \texttt{fin}.
            \item Subtracts divisor $\times$ digit from \texttt{comp} and updates it.
        \end{itemize}
    \item If both input numbers were negative, I divided their absolute values.
    \item If only one of the inputs was negative, I divided their absolute values and prefixed the result with a minus sign.
    \item Finally, the result was passed through a function to remove any leading zeros before returning it.
    \item I wrapped the result into a new \texttt{AInteger} and returns it.
\end{itemize}

\section*{AFloat.java / Float class}
I created the file \texttt{AFloat.java} inside the \texttt{arbitraryarithmetic} directory. To include it in the same package, I added the statement \texttt{package arbitraryarithmetic;} at the top of the file.

\subsection*{Properties or Members}
It has a protected member named \texttt{'num'} of type String.

\subsection*{Constructors}
\begin{itemize}
    \item Default constructor \texttt{AFloat()} that initializes the instance with value 0.0.
    \item Constructor \texttt{AFloat(String num)} that initializes the instance by the number whose string representation is given by 'num'.
    \item Copy constructor that creates an instance of \texttt{AFloat}.
\end{itemize}

\subsection*{Methods / Functions}
\begin{description}
    \item[\texttt{parse(String num)}] A static function that returns an instance of \texttt{AFloat} class.
    \item[\texttt{get()}] Returns the member of the object.
    \item[\texttt{set(String num)}] Sets the member of the object to \texttt{num}.
    \item[\texttt{decimal\_digits()}] Returns the number of decimal digits of their member.
    \item[\texttt{decimal\_point()}] Confirms whether the member has decimal point.
    \item[\texttt{remove\_extra\_zeroes(String str)}] Function which removes extra zeroes at the end of decimal places.
    \item[\texttt{add(AFloat other)}] Adds any two floating-point numbers of arbitrary length.
    \item[\texttt{sub(AFloat other)}] Subtracts any two floating-point numbers of arbitrary length.
    \item[\texttt{mul(AFloat other)}] Multiplies any two floating-point numbers of arbitrary length.
    \item[\texttt{div(AFloat other)}] Divides any two floating-point numbers of arbitrary length.
    \item[\texttt{truncate(String str)}] Truncates the result to 30 decimal places.
\end{description}

\subsection*{A deeper look into methods}
\subsubsection*{\texttt{add(AFloat other)}}
\begin{itemize}
    \item It implements the \texttt{add()} method in \texttt{AFloat.java} to perform addition of two floating-point numbers represented as strings.
    \item First, It ensures both numbers are in a proper floating-point format
    \item If a number didn't already contain a decimal point, It appends one to ensure consistent processing.
    \item It calculates the decimal lengths of both operands using the \texttt{decimal\_digits()} method.
    \item If both operands had the same number of digits after the decimal:
        \begin{itemize}
            \item It removes the decimal points and adds the integer parts of the numbers.
            \item The result is then formatted by inserting the decimal point at the correct position based on the original decimal length.
        \end{itemize}
    \item If operand 1 has a greater decimal length than operand 2:
        \begin{itemize}
            \item It pads operand 2 with zeros to match the decimal length of operand 1.
            \item The result is calculated in the same way as the equal decimal length case.
        \end{itemize}
    \item If operand 2 has a greater decimal length than operand 1:
        \begin{itemize}
            \item It pads operand 1 with zeros to match the decimal length of operand 2.
            \item The result is calculated in the same way as the equal decimal length case.
        \end{itemize}
    \item After the addition of the integer parts, it handles the following:
        \begin{itemize}
            \item If the result has more digits than the original decimal length, it inserts the decimal point at the proper position.
            \item If the result has fewer digits, it adds the necessary number of leading zeros before the decimal point.
            \item If the result is negative, the negative sign is handled and inserted properly.
        \end{itemize}
    \item After constructing the result string, It removes any unnecessary leading zeroes and uses the \texttt{remove\_extra\_zeroes()} and \texttt{truncate()} functions to clean up the result before returning it.
    \item Finally, it wraps the result into a new \texttt{AFloat} object and returns it.
\end{itemize}

\subsubsection*{\texttt{sub(AFloat other)}}
\begin{itemize}
    \item Convert both operands to strings and ensure they have a decimal point.
    \item If either operand doesn't have a decimal point, it adds one.
    \item It calculates the decimal lengths of both operands using the \texttt{decimal\_digits()} method.
    \item If both operands have equal decimal lengths:
        \begin{itemize}
            \item It removes the decimal points and performs subtraction by treating the numbers as integers.
            \item The result is then formatted, ensuring the decimal point is inserted correctly based on the original decimal length.
            \item If the result has more digits than the decimal length, the decimal point is inserted at the correct position.
            \item If the result is negative, the negative sign is handled properly and inserted in the correct place.
            \item If the result is smaller than the decimal length, zeroes are padded as required.
        \end{itemize}
    \item If the decimal lengths are not equal:
        \begin{itemize}
            \item The operand with the smaller decimal length is padded with zeros to make both lengths equal.
            \item Subtraction is then performed as before, with proper formatting of the result, including zero padding where necessary.
        \end{itemize}
    \item It removes unnecessary leading zeros.
    \item It ensures no trailing zeros remain using the \texttt{remove\_extra\_zeroes()} method.
    \item It returns the final result as a new \texttt{AFloat} object.
\end{itemize}

\subsubsection*{\texttt{mul(AFloat other)}}
\begin{itemize}
    \item Convert both operands to strings and ensure they contain a decimal point.
    \item If either operand doesn't have a decimal point, it adds one.
    \item It calculates the decimal lengths of both operands using the \texttt{decimal\_digits()} method
    \item The decimal point is removed from both operands by converting them into integer-like strings.
    \item Leading zeros are also removed from both operands before multiplication.
    \item The operands are converted into \texttt{AInteger} objects and multiplied using \texttt{AInteger:: mul(AInteger odd)}.
    \item The result is stored and processed based on the total number of decimal places:
        \begin{itemize}
            \item The total decimal length is the sum of the decimal lengths of both operands.
            \item If the result has more digits than the total decimal length, the decimal point is inserted in the correct position.
            \item If the result has fewer digits than the total decimal length, it pads the result with leading zeros and places the decimal point correctly.
        \end{itemize}
    \item If the result is negative, the negative sign is handled properly and inserted in the correct place, with additional formatting for the decimal point.
    \item It removes unnecessary leading zeros.
    \item It ensures no trailing zeros remain using the \texttt{remove\_extra\_zeroes()} method.
    \item The result is truncated using the \texttt{truncate()} method to ensure the correct format.
    \item Finally, it wraps the result into a new \texttt{AFloat} object and returns it.
\end{itemize}

\subsection*{\texttt{div(AFloat other)}}
\begin{itemize}
    \item First, it ensures both operands have a decimal point.
    \item If either operand doesn't have a decimal point, \texttt{.0} is appended to it.
    \item The operands are parsed into \texttt{AFloat} objects to determine their decimal lengths using the \texttt{decimal\_digits()} method.
    \item The decimal points are removed by converting both operands into integer-like strings.
    \item Leading zeros are removed from both operands before performing division.
    \item The division operation is performed by using \texttt{AInteger :: div (AInteger other)} for both the quotient and the remainder.
    \item If division by zero is attempted, an \texttt{IllegalArgumentException} is thrown.
    \item If the first operand is zero, the result is zero.
    \item The sign of the result is determined by the signs of the operands:
        \begin{itemize}
            \item If one operand is negative, the result will be negative.
            \item If both operands are negative, the result will be positive.
        \end{itemize}
    \item The remainder is computed and processed to remove any negative sign and ensure it is handled correctly during division.
    \item The long division method is implemented to compute the result up to 30 decimal places:
        \begin{itemize}
            \item The remainder is repeatedly divided by the divisor, and the quotient is appended to the decimal part of the result.
            \item If the remainder becomes zero, the loop terminates.
        \end{itemize}
    \item he result is formatted correctly:
        \begin{itemize}
            \item The quotient and the decimal part are concatenated.
            \item If there are extra zeros, they are added or removed as necessary, and the decimal point is inserted in the correct position.
        \end{itemize}
    \item If the result is negative, the minus sign is handled correctly, and the number is formatted accordingly.
    \item Leading zeros are removed from the result using a custom method to ensure the correct formatting.
    \item After the result is calculated, it is verified by multiplying it back with the divisor to ensure no extra zeros are present. If necessary, extra zeros are removed.
    \item The final result is returned as an \texttt{AFloat} object.
\end{itemize}

\section*{MyInfArith.java}
\begin{itemize}
    \item It is a Java file that imports the package and runs test cases.
    \item It ensures that the input passed is only an integer or a float.
    \item It takes command-line arguments and performs arithmetic operations.
\end{itemize}

\section*{my\_exe}
\begin{itemize}
    \item It is a Python script used to run test cases.
    \item It uses \texttt{MyInfArith.java} to run the test cases.
\end{itemize}

\section*{build.xml}
\begin{itemize}
    \item Defines a project named \textbf{"Arbitrary Arithmetic Operations"} with the default target \texttt{jar}.
    \item Sets both the source and build directories to the \textbf{current folder}, so \texttt{.class} files go alongside \texttt{.java} files.
    \item The \texttt{compile} target compiles \texttt{arbitraryarithmetic/*.java} into the same directory.
    \item The \texttt{jar} target creates \texttt{aarithmetic.jar} containing all \texttt{.class} files from the current directory tree.
    \item The \texttt{clean} target deletes the \texttt{.jar} file and all \texttt{.class} files from the directory.
\end{itemize}

\section*{aarithmetic.jar}
The \texttt{aarithmetic.jar} library can be linked to any \textbf{executable/library}.


\end{document}
